% Source: https://artofproblemsolving.com/community/c3h1503979_tessellation_trouble
% Stable identity: aops-topic-1503979-q01
The interior angle of a regular $n$-gon is
\[
\alpha=\dfrac{(n-2)\cdot180^\circ}{n}=180^\circ-\dfrac{360^\circ}{n},
\]
so $60^\circ\le\alpha<180^\circ$.

\textbf{Necessary condition.} Take a vertex $P$ of one tile. The tiles around $P$ fill the full $360^\circ$ there. Each of them either has a vertex at $P$, contributing an angle $\alpha$, or has $P$ inside one of its edges, contributing $180^\circ$. There can be at most one tile of the second kind, since two such tiles would leave no room for the tile with a vertex at $P$. So
\[
j\alpha=360^\circ\qquad\text{or}\qquad j\alpha=180^\circ
\]
for some positive integer $j$.
\begin{itemize}
\item In the first case $j=\dfrac{360}{\alpha}=\dfrac{2n}{n-2}=2+\dfrac{4}{n-2}$. So $n-2$ divides $4$, and $n\in\{3,4,6\}$.
\item In the second case $j=\dfrac{n}{n-2}=1+\dfrac{2}{n-2}$. So $n-2$ divides $2$, and $n\in\{3,4\}$.
\end{itemize}

\textbf{These values work.} Equilateral triangles ($6$ at a vertex), squares ($4$ at a vertex) and regular hexagons ($3$ at a vertex) tile the plane:

\begin{asy}
usepackage("amsmath");
size(8cm);
defaultpen(fontsize(10pt));
pair P1=(0,0);
for(int k=0;k<6;++k){ draw(P1--P1+dir(60k),linewidth(0.8)); draw(P1+dir(60k)--P1+dir(60k+60),linewidth(0.8)); }
draw(arc(P1,0.22,0,60),red);
label("$6\times60^\circ$",(0,-1.35));
pair P2=(3.1,0);
draw(shift(P2)*((-0.9,-0.9)--(0.9,-0.9)--(0.9,0.9)--(-0.9,0.9)--cycle),linewidth(0.8));
draw((P2.x-0.9,0)--(P2.x+0.9,0),linewidth(0.8)); draw((P2.x,-0.9)--(P2.x,0.9),linewidth(0.8));
draw(arc(P2,0.22,0,90),red);
label("$4\times90^\circ$",(P2.x,-1.35));
pair P3=(6.2,-0.05); real s=0.62;
for(int t=0;t<3;++t){ pair c=P3+s*dir(90+120t); path h=c+s*dir(270+120t);
  for(int k=1;k<6;++k) h=h--(c+s*dir(270+120t+60k));
  draw(h--cycle,linewidth(0.8)); }
draw(arc(P3,0.22,-30,90),red);
label("$3\times120^\circ$",(P3.x,-1.35));
dot(P1);dot(P2);dot(P3);
label("\mbox{}",(min(currentpicture,user=true)+max(currentpicture,user=true))/2,black); // keeps the text after the figure black when LaTeX typesets the labels inline
\end{asy}

Hence the plane can be tiled by a regular $n$-gon exactly for $\boxed{n=3,\ 4,\ 6}$.
